% ======================================================================
% Appendix C — Self-contained verification core
% ======================================================================

\section{A self-contained verification core}\label{sec:self}

This appendix lets a reader with only this paper verify the sharpest
negative claims without the companion monograph. It collects the
elementary identity behind the four forms, the runner formulation with
its certificate, the explicit statement of the sum-only equality, and
one-line hand checks of the headline witnesses.

\subsection{The identity and the runner formulation}

\begin{lemma}[depth identity]\label{lem:depthid}
For every $t\in\T$,
$\max_{1\le j\le n}\wnc{\tfrac12-v_jt}=\tfrac12-\min_{1\le j\le n}\wnc{v_jt}$.
Consequently $\delta(V)=\tfrac12-\lambda(V)$ with
$\delta(V)=\operatorname{dist}_\infty\!\bigl(\R v,\;\Z^n+\tfrac12\mathbf{1}
\bigr)$, and the $\delta$-form of \S\ref{sec:prelim} is exactly
$\lrc_n$.
\end{lemma}

\begin{proof}
Each map $x\mapsto\wnc{x}$ satisfies
$\wnc{\tfrac12-x}=\tfrac12-\wnc{x}$: reflecting about $\tfrac12$
exchanges the two half-intervals and fixes the midpoint. Taking $x=v_jt$
and maximizing over $j$ gives the displayed identity (a maximum of
quantities each equal to $\tfrac12$ minus the corresponding term is
$\tfrac12$ minus the minimum). The distance form is the same statement
along the trajectory line $tv$, $t\in\R$.
\end{proof}

The computations of this paper use the \emph{runner formulation}
\eqref{eq:runnerform}, restated here:
\[
D(c)\;=\;\min_{\sigma\in[0,1)}\;\max\Bigl(\wnc{\sigma},\;
\max_{2\le j\le n}\wnc{c_j-v_j\sigma}\Bigr).
\]
Two facts make it auditable. First, it is the quotient-norm depth
$\operatorname{dist}_{\mathcal{X}_v}(c,\Lambda)$ written in the
$c$-coordinates of \S\ref{sec:prelim}: the linear forms $\phi^{1j}$,
$\phi^{ij}$ listed there generate the norm, the parameter $\sigma$ is the
free coordinate $x_1$ of a lift, and minimizing over $\sigma$ minimizes
over the lift. Second, it is \emph{Lipschitz-auditable}: each term
$\wnc{c_j-v_j\sigma}$ is $v_j$-Lipschitz in $\sigma$ and the whole
objective is piecewise linear with breakpoints at the half-integers
$c_j-\tfrac12-v_j\sigma\in\Z$, so the minimum is attained at a rational
breakpoint or crossing and is exactly computable by hand or by exact
rational arithmetic. The independent linear-form/lattice-search
implementation of the program agrees with it exactly on every instance
used in this paper (forty random rational points per instance).

\subsection{The sum-only equality, explicitly}

\begin{theorem}[sum-only equality; explicit form]\label{thm:equality-explicit}
Let $V=(v_1,\dots,v_n)$ have distinct positive entries with $\gcd(V)=1$
and $0<M(V)<\tfrac12$. For each unordered pair $\{u,w\}\subset[n]$ write
$N_{uw}=v_u+v_w$ and
\[
\tau(u,w)\;=\;\max_{0\le k<N_{uw}}\;\min_{p\in[n]\setminus\{w\}}\;
\wnc{\frac{v_pk}{N_{uw}}}.
\]
Then $M(V)=\max_{\{u,w\}}\tau(u,w)$, and consequently $\lrc_n$ holds
for all $V$ if and only if for every $V$ there is a pair $\{u,w\}$ with
$\tau(u,w)\ge 1/(n+1)$.
\end{theorem}

The two ingredients (proved in the companion monograph, Theorem~1.1
there; the finitization is what every instance of \S\ref{sec:tm-covering}
entered through) are: (i) at a global optimum $t^\ast$ with
$0<M(V)<\tfrac12$ the binding runners can be chosen two-sided---one
rising, one falling---which forces $(v_u+v_w)t^\ast\in\Z$ for that pair,
so $t^\ast$ lies on the \emph{sum} lattice of denominator $v_u+v_w$; and
(ii) on that lattice the two binding runners are exact reflections
($\wnc{v_uk/N}=\wnc{v_wk/N}$ for all $k$, since $v_u\equiv-v_w \pmod N$),
so the $n$-runner minimum equals an $(n-1)$-runner minimum. The
companion's Lemma~2.1 also records the boundary cases $M(V)=\tfrac12$
(all speeds odd) and $M(V)=0$ (impossible).

\subsection{Hand checks of the headline witnesses}

All the strict witnesses of Table~\ref{tab:refute} are one-minute hand
computations with \eqref{eq:runnerform}; the two closed values need the
branch-and-bound half of their certificates as well
(\S\ref{sec:bb}), but the witness halves are hand checks too.

\emph{The harmonic argmax value $16/47$.} Take $v=(1,2,3,4,5)$ and
$c=(\tfrac{122}{1175},\tfrac{688}{1175},\tfrac{284}{1175},
\tfrac{80}{1175})$. At $\sigma=\tfrac{96}{1175}$ (working in
$1175$-ths, with $c=(\tfrac{122}{1175},\tfrac{688}{1175},
\tfrac{284}{1175},\tfrac{80}{1175})$):
$\wnc\sigma=\tfrac{96}{1175}$;
$\wnc{c_2-2\sigma}=\wnc{\tfrac{122-192}{1175}}=\tfrac{70}{1175}$;
$\wnc{c_3-3\sigma}=\wnc{\tfrac{688-288}{1175}}=\tfrac{400}{1175}$;
$\wnc{c_4-4\sigma}=\wnc{\tfrac{284-384}{1175}}=\tfrac{100}{1175}$;
$\wnc{c_5-5\sigma}=\wnc{\tfrac{80-480}{1175}}=\tfrac{400}{1175}$.
The maximum is $\tfrac{400}{1175}=\tfrac{16}{47}$, attained twice
(runners $3$ and $5$, slopes $+3$ and $-5$ opposite), and it is
minimized: any local shift of $\sigma$ increases one of the two. Hence
$D(c)=16/47>\tfrac13=\mathrm{thr}_5$. The pair-form side of the same
point is the quadruple binding of \S\ref{sec:weights}:
$\phi^{45}(c-z)=-16/47$ at $z=(0,0,0,-1)$, $\phi^{34}(c-z)=-16/47$ at
$z=(0,1,0,0)$, $\phi^{35}(c-z)=+16/47$ at $z=0$, and
$\phi^{25}(c-z)=+16/47$ at $z=(0,1,1,1)$, each a two-term integer
evaluation.

\emph{The rung-2 argmax value $17/47$.} Take
$v=(1,2,3,4,5,12)$ and
$c=(\tfrac{4714}{5875},\tfrac18,\tfrac58,\tfrac6{25},
\tfrac{2034}{5875})$. At $\sigma=\tfrac{482}{5875}$ the six terms are
$\tfrac{482}{5875}$, $\tfrac{2125}{5875}$, $\tfrac{5693}{47000}$,
$\tfrac{13951}{47000}$, $\tfrac{1000}{5875}$, $\tfrac{2125}{5875}$:
the maximum $\tfrac{2125}{5875}=\tfrac{17}{47}$ is attained twice
(runners $2$ and $12$, slopes $+2$ and $-12$ opposite) and is
minimized, so $D(c)=17/47>\tfrac5{14}=\mathrm{thr}_6$. The binding
triple $\{2,5,12\}$ and its translates are recorded in
\S\ref{sec:weights}.

\emph{The strict witnesses.} For $v=(1,2,3,4,5,6)$ take
$c=(0,\tfrac5{16},\tfrac{15}{16},\tfrac34,\tfrac12)$ and
$\sigma=\tfrac{23}{160}$: the six terms are
$\tfrac{23}{160}$, $\tfrac{23}{80}$, $\tfrac{19}{160}$, $\tfrac{29}{80}$,
$\tfrac1{32}$, $\tfrac{29}{80}$, so $D(c)=29/80>\tfrac5{14}$ (runners
$4$ and $6$ tie at the maximum with opposite slopes). For
$v=(1,2,3,4,5,7)$ take
$c=(\tfrac1{16},\tfrac9{16},0,\tfrac38,0)$ and $\sigma=\tfrac1{11}$:
the six terms are $\tfrac1{11}$, $\tfrac{21}{176}$, $\tfrac{51}{176}$,
$\tfrac4{11}$, $\tfrac7{88}$, $\tfrac4{11}$, so $D(c)=4/11>
\tfrac5{14}$. For $v=(1,2,3,4,5,9)$ take
$c=(\tfrac1{32},\tfrac{15}{32},\tfrac{1383}{10000},
\tfrac{32921}{40000},\tfrac{95473}{200000})$ and
$\sigma=\tfrac{18577}{200000}$: the six terms are
$\tfrac{18577}{200000}$, $\tfrac{3863}{25000}$,
$\tfrac{38019}{200000}$, $\tfrac{5831}{25000}$, $\tfrac{1793}{5000}$,
$\tfrac{1793}{5000}$, so $D(c)=1793/5000>\tfrac5{14}$ (runners $5$ and
$9$ tie at the maximum).

\emph{The $7/22$ center class (the $m=10$ ladder closure's witness).}
Take $v=(1,2,3,4,10)$ and
$c=(\tfrac12,0,\tfrac12,\tfrac12)$. At $\sigma=\tfrac{7}{22}$ the five
terms of \eqref{eq:runnerform} are $\tfrac{7}{22}$, $\tfrac{3}{22}$,
$\tfrac{1}{22}$, $\tfrac{5}{22}$, $\tfrac{7}{22}$: explicitly
$\wnc\sigma=\tfrac7{22}$; $\wnc{\tfrac12-2\sigma}=\tfrac3{22}$;
$\wnc{-3\sigma}=\tfrac1{22}$;
$\wnc{\tfrac12-4\sigma}=\tfrac5{22}$; and
$\tfrac12-10\sigma=-\tfrac{59}{22}\equiv\tfrac7{22}\pmod 1$, so the
fifth term is $\tfrac7{22}$. The maximum is $\tfrac7{22}$, and it is
minimized: the binding terms are the $\sigma$-term (slope $+1$) and the
fifth term (slope $-10$), opposite signs, so any local shift of
$\sigma$ increases one of them. Hence $D(c)=7/22$, and the
branch-and-bound certificate of \S\ref{sec:bb} ($\rho\le7/22$, the
tree emptying) upgrades this to $\rho(1,2,3,4,10)=\tfrac{7}{22}$
exactly; since $\delta=1/2-\lambda=1/2-2/11=7/22$, the deepest hole at
$m=10$ is the center class itself. The pair-form side is the bisection
of \S\ref{sec:weights}: the pair $(1,10)$ binds at two translates with
opposite signs, $2\cdot 11\cdot\rho=7\in\Z$.

\emph{The balance-family witness at the harmonic argmax.} Take
$c=c(35)=(\tfrac{46}{47},\tfrac{92}{235},\tfrac{231}{235},
\tfrac{35}{47})$. At $\sigma=\tfrac{4}{235}$ the five terms of
\eqref{eq:runnerform} are $\tfrac{4}{235}$, $\tfrac{13}{235}$,
$\tfrac{80}{235}$, $\tfrac{20}{235}$, $\tfrac{80}{235}$: explicitly
$\wnc\sigma=\tfrac4{235}$;
$\wnc{\tfrac{46}{47}-2\sigma}=\wnc{\tfrac{222}{235}}=\tfrac{13}{235}$;
$\wnc{\tfrac{92}{235}-3\sigma}=\tfrac{80}{235}$;
$\wnc{\tfrac{231}{235}-4\sigma}=\wnc{\tfrac{215}{235}}=\tfrac{20}{235}$;
and $\wnc{\tfrac{35}{47}-5\sigma}=\wnc{\tfrac{155}{235}}=\tfrac{80}{235}$.
The maximum is $\tfrac{80}{235}=\tfrac{16}{47}$, attained with the
third and fifth terms binding (slopes $-3$ and $+5$, opposite), and the
tie over all three dips of the triple $\{3,4,5\}$---forced by the
identity along the family---gives the same value at every member; the
two-sided certificate of \S\ref{sec:bb} then makes $16/47$ the exact
covering radius.

Every value in
Table~\ref{tab:refute} was additionally re-asserted by two independent
implementations in the provenance run (Appendix~\ref{sec:prov}), and
the closure runs' witnesses by the cross-validation run described
there.
